\begingroup
\small
\setlength{\tabcolsep}{4pt}
\renewcommand{\arraystretch}{1.16}
\begin{longtable}{
>{\raggedright\arraybackslash}p{0.25\textwidth}
>{\raggedright\arraybackslash}p{0.44\textwidth}
>{\raggedright\arraybackslash}p{0.23\textwidth}}
\caption{Dostępne kolumny w surowych zbiorach danych oraz sposób ich parsowania}
\label{tab:raw-dataset-schema-inventory}\\
\hline
\textbf{Zbiór} & \textbf{Dostępne kolumny} & \textbf{Sposób parsowania} \\
\hline
\endfirsthead
\hline
\textbf{Zbiór} & \textbf{Dostępne kolumny} & \textbf{Sposób parsowania} \\
\hline
\endhead
\hline
\endfoot
\texttt{enron} &
\texttt{file}, \texttt{message} &
Wydzielenie tematu i treści z pola \texttt{message}. \\

\texttt{trec\_2007} &
\texttt{label}, \texttt{origin} &
Wydzielenie tematu i treści z pola \texttt{origin}. \\

\texttt{ceas\_2008} &
\texttt{sender}, \texttt{receiver}, \texttt{date}, \texttt{subject}, \texttt{body}, \texttt{label}, \texttt{urls} &
Bezpośrednie mapowanie pól \texttt{subject}, \texttt{body}, \texttt{label}. \\

\begin{tabular}[t]{@{}l@{}}\texttt{spam}\texttt{\_assassin}\end{tabular} &
\texttt{text}, \texttt{target} &
Wydzielenie tematu i treści z pola \texttt{text}. \\

\texttt{spam\_ham} &
\texttt{label}, \texttt{text}, \texttt{label\_num} &
Wydzielenie tematu z prefiksu \texttt{Subject:}. \\

\begin{tabular}[t]{@{}l@{}}\texttt{fraudulent}\texttt{\_email}\\\texttt{\_corpus}\end{tabular} &
wiadomości zapisane kolejno z nagłówkami &
Wydzielenie tematu z nagłówków i treści z wiadomości. \\

\texttt{nazario} &
\texttt{sender}, \texttt{receiver}, \texttt{date}, \texttt{subject}, \texttt{body}, \texttt{label}, \texttt{urls} &
Bezpośrednie mapowanie pól \texttt{subject}, \texttt{body}, \texttt{label}. \\

\texttt{ling} &
\texttt{subject}, \texttt{message}, \texttt{label} &
Przeniesienie \texttt{message} do pola \texttt{body} \\
\end{longtable}
\endgroup
